\section{Representation Learning: Contrastive y Predictive Objectives}

Una vez que se tienen code discretos, el segundo artículo compara dos rutas que no dependen de etiquetas. El contrastive learning acerca entre sí distintas augmentation de un mismo sonido y aleja los sonidos diferentes; el generative/predictive learning, en cambio, exige que el latent code prediga el futuro. Ambos buscan aprender una representation transferible, pero conservan información distinta: el primero codifica la invariancia que se define, y el segundo codifica el estado útil para que el futuro sea predecible.

\subsection{Qué acerca el Contrastive Learning}

Para un anchor $x$, una positive augmentation $x^+$ y un conjunto de negatives $x_j^-$, puede escribirse un InfoNCE-style objective:
\[
\mathcal{L}_{\mathrm{NCE}}
=-\log
\frac{\exp(\mathrm{sim}(f(x),f(x^+))/\tau)}
{\exp(\mathrm{sim}(f(x),f(x^+))/\tau)
+\sum_j\exp(\mathrm{sim}(f(x),f(x_j^-))/\tau)}.
\]
Aquí $f$ es el encoder, $\mathrm{sim}$ es la similitud y $\tau$ es la temperature. La augmentation no es una operación neutral: si un pitch shift se trata como positive, la representación ignorará más la altura tonal; esto puede favorecer la scene classification, pero perjudicar la pitch estimation.

\lecturefigure{slide-32-contrastive-architecture.jpg}{Contrastive audio representation: augmentation pairs, encoder y objetivo de similitud.}{Video oficial de Stanford Online 00:30:25--00:32:25.}

\begin{warningbox}{Invariant to what?}
El contrastive objective no descubre automáticamente la «semántica correcta»: aprende la invariancia que define la augmentation. En audio, si pitch, tempo, timbre, noise y time shift deben ignorarse depende de la tarea posterior; una augmentation equivocada elimina activamente información esencial.
\end{warningbox}

\subsection{La estructura común de Generative, Clustering y Speech Models}

En clase se colocan Jukebox, VQ-VAE, k-means, wav2vec-like codes, Tacotron y los generative speech models en un mismo diagrama comparativo. Aunque son muchos nombres, el problema puede reducirse a tres columnas: quién produce las latent units, quién modela las units y quién se encarga de la reconstrucción o la lectura. Solo así se ve que «el Transformer predice code» y «el Transformer produce directamente la waveform» se sitúan en niveles de abstracción distintos.

\lecturefigure{slide-33-audio-methods-comparison.jpg}{Audio methods comparison: cómo se reparten el trabajo los distintos sistemas entre encoder, discretizer, sequence model y decoder.}{Video oficial de Stanford Online 00:31:58--00:33:21.}

\begin{table}[H]
\centering
\small
\caption{Comparación de mecanismos de las rutas autosupervisadas/generativas para audio.}
\begin{tabularx}{\textwidth}{L{0.20\textwidth}L{0.22\textwidth}L{0.24\textwidth}X}
\toprule
Ruta & Señal de aprendizaje & Contenido que principalmente conserva & Riesgo principal \\
\midrule
Contrastive & Distinguir el positive de los negatives & Atributos que permanecen estables tras la augmentation & False negatives, invariance equivocada, dependencia del lote. \\
Next-code prediction & Predecir los codes discretos futuros & Estados con poder predictivo sobre el futuro & Tokenizer bottleneck, exposure bias, atajos locales. \\
Waveform prediction & Predecir el siguiente sample & Detalle a nivel de muestra y dinámica local & Secuencias extremadamente largas; desajuste entre la metric y la percepción auditiva. \\
Reconstruction & Recuperar la entrada a partir del latent & Lo que el decoder puede reconstruir & Tiende a conservar detalles de bajo nivel e ignorar la semántica de la tarea. \\
\bottomrule
\end{tabularx}
\end{table}

\lecturefigure{slide-34-textless-nlp.jpg}{Textless NLP / generative speech: modelado de secuencias con speech units discretas y su posterior decodificación.}{Video oficial de Stanford Online 00:33:22--00:33:41.}

\subsection{Por qué predecir el futuro puede producir un Summary}

Si el latent state actual puede predecir con exactitud los codes futuros, debe comprimir la información de la historia que es relevante para lo que sigue. El ponente lo llama «better prediction, better summarization». Formalmente, se busca que la representation $h_t$ cumpla
\[
p(z_{t+1:t+K}\mid z_{\le t})
\approx p(z_{t+1:t+K}\mid h_t),
\]
donde $z_t$ es un code discreto. Esta aproximación expresa el objetivo del predictive state; no garantiza que $h_t$ sea suficiente para todas las tareas futuras, y también puede ocurrir que solo aprenda patrones predecibles de corto plazo presentes en los datos.

\lecturefigure{slide-35-long-dependency-cluster-model.jpg}{Aprender long-term dependency con una cluster-code sequence: la capacidad predictiva como indicador sustituto de la representation quality.}{Video oficial de Stanford Online 00:33:42--00:34:51.}

\teachervoice{«Cuanto mejor la predicción, mejor el summary» es un supuesto central del curso, no una definición. Si el futuro depende sobre todo de la textura local, el modelo puede obtener un buen prediction loss sin comprender la semántica de la escena.}

\subsection{Cómo comparar con equidad las Learned Representations}

El artículo procura fijar el dataset y la encoder family, y entrena una linear head sobre la representación congelada. Si la representación es $h=f_\theta(x)$, el linear probe es
\[
\hat{y}=\operatorname{softmax}(Wh+b),
\]
y durante el entrenamiento solo se actualizan $W,b$. Si la accuracy es alta, la información de las etiquetas puede leerse linealmente; si es baja, puede deberse a que falta información o simplemente a que se requiere un nonlinear readout. Fijar el encoder y los datos reduce el confounding, pero el objective de preentrenamiento, la augmentation, la hyperparameter search y el code rate siguen influyendo en los resultados.

\lecturefigure{slide-36-unsupervised-representation-evidence.jpg}{Diseño de la controlled comparison: same encoder, same dataset, linear readout.}{Video oficial de Stanford Online 00:34:52--00:36:07.}

\lecturefigure{slide-37-same-encoder-same-data-surface.jpg}{Resultados de la representación: supervised 71.3\%, contrastive 55.6\%, next-code Transformer 60.1\%.}{Video oficial de Stanford Online 00:36:08--00:37:09; arXiv:2010.11459.}

\begin{knowledgebox}{Lectura de la figura: primero el control, después el gap}
Usar el mismo encoder y el mismo dataset hace que el objective sea la principal fuente de variación. El next-code Transformer supera a ese contrastive setup, pero sigue unos 11.2 points por debajo del supervised baseline. La surface de datos de entrenamiento/hidden-size de la derecha muestra que la escala modifica el error, pero no demuestra directamente que «más datos sin etiquetar eliminan necesariamente el gap»; esa es la scaling hypothesis que plantea el ponente.
\end{knowledgebox}

\begin{warningbox}{Los límites de las conclusiones del linear probe}
71.3\%, 55.6\% y 60.1\% son la accuracy para el dataset/encoder/probe actual. No comparan generation quality, few-shot transfer, robustness, calibration ni compute efficiency, y tampoco permiten afirmar que la supervised representation y la unsupervised representation difieran solo en un porcentaje fijo.
\end{warningbox}

\teachervoice{Verma sitúa explícitamente «más datos reducirán el gap del 10\%» después de los resultados, como una conjetura. Estas notas conservan ese orden: primero se reporta la tabla controlada y después se discute el scaling, en lugar de presentar la conjetura como una conclusión ya verificada por el artículo.}

\subsection{Resumen de la sección}

El contrastive learning define la invariance mediante la augmentation; el predictive learning define el estado suficiente mediante los future codes. El VQ tokenizer, el objective y el probe determinan juntos el resultado; el control con los mismos datos y el mismo encoder es importante, pero aun así no permite extrapolar la linear accuracy a una representation quality general.
